\doublespacing
\section{Culture-talk and status-based stratification}
The empirical issue that is the focus of this paper pertains to the ``behavioral implications'' of the aesthetic consumption of popular (and fine arts) culture \citep{fine77}.  This phenomenon can best be characterized through the frequently noted observation that cultural knowledge obtained from aesthetic consumption experiences will become a likely content used for the initiation and maintenance of conversational interaction \citep{dimaggio87, collins75}.  I refer to this (presumably) pervasive phenomenon as \emph{culture-talk}.  Culture-talk can be defined as the deployment of cultural knowledge obtained from aesthetic consumption experiences as a resource to generate conversation.  In a nutshell:  the discussion of ``cultural matters'' with friends and acquaintances.  

Culture-talk is one of the most common facilitators of \emph{sociability} in late-modern societies.  Sociability is social interaction that is not driven by explicit instrumental purposes, but that is keyed to the continuing repetition of episodes of social intercourse for their own sake \citep{simmel49}. This is a form of social intercourse explicitly dissociated from direct instrumental pursuits, as a \emph{form} of association ``which does not have a strictly economic or business purpose'' but access to which (and restrictions on) are a key marker of the status group boundaries \citep[114]{weber94}. It is in this respect that sociability can be thought of as the ``play form''  \citep{simmel49} of the more general process of association (coming together from some purpose).  This type of interaction would be difficult to successfully realize by relying on traditional ``instrumental'' contents (e.g. work, politics, etc.), but requires instead ``ludic'' content that serves to simply sustain the focus of attention in the interaction itself \citep{collins04}.  I argue that the ``entertainment talk'' \citep[124-125]{collins75} which makes up most of the conversational matters of the more cosmopolitan middle and upper classes is one of the forms of cultural knowledge that is most apt for these ``playful'' purposes. 

\subsection{The interactional functions of culture-talk}
As already intimated, culture-talk as a interaction-based mechanism for the creation of status group boundaries is at the center of Collins' microinteractionist theory of culture and stratification \citep{collins75, collins81, collins00, collins04}. The key claim of this line of theory is that sociability is unequally distributed by status, because it is premised on unequal access to and unequal opportunity to accumulate the relevant set of dispositions required to translate consumption experiences into conversation for its own sake. According to Colllins, status-differentiated sociability is primarily built around what he refers to as ``conversation for enjoyment'' \citep[124]{collins75}.  These types of conversational rituals are important because they signal a temporary removal from instrumental realms laden with power-dynamics, status differences and ``arms-length'' relationships built for explicit, narrow purposes.  Instead, conversation for enjoyment indicate that we have entered into a homo-social realm based on (situational) status equality and the privilege of status inclusion \citep{milner04}.  For this reason, it is to be expected  that culture-based talk should be more likely to happen among friends than among strangers \citep[124]{collins75}. This makes culture-talk a primary ingredient in the constitution and reproduction of status-segregated micro-orders. In this respect, the often noted phenomenon of ``homophily'' is at least partially constituted by ``choice'' processes based on shared tastes as it is by structural processes that induce propinquity and bias interaction \citep{lazarsfeld_merton54}\footnote{Nevertheless, certain forms of culture-talk---especially those involving cultural content presumed to be spread evenly across groups---provide the means to provisionally relax status groups boundaries by providing a baseline form of cultural knowledge that can allow otherwise socially distant persons to ``break the ice'' \citep{erickson96}. This makes social structures a bit more ``elastic'' at least for the duration of the encounter \citep[85]{goffman66}.} 

Status-based exclusion organized around conversational content occurs largely through micro-mechanisms governing entry (or determining systematic exclusion from) class-structured micro-worlds organized around specific conversational interactions \citep{collins75, collins00}; lack of familiarity with the relevant conversational currency or inability to produce discourse of the appropriate sort results in exclusion from the relevant circle \citep{erickson96, purcell07}.  In this respect, to the extent that the experiential and phenomenological ``stuff'' of stratification is essentially based on talk then we should expect much of this talk to be about aesthetic experiences.  This is the reason why we can consider culture-talk to be ``\ldots a form of ritual identification and a means of construction social relations\ldots It helps to establish networks of trusting relations that facilitate group mobilization and the attainment of such social rewards as desirable spouses and prestigious jobs'' \citep[443]{dimaggio87}.  Culture-talk thus serves to increase the cohesion and social closure of small groups, providing a way to situationally segment them from the larger social structure.

Thus, according to this line of analysis, routine engagement in culture-talk serves as the primary mechanism responsible for the generation and reproduction of micro-solidarities that hold status-groups together \citep{fine77, ostrower98}.  The primary (unintended) function of aesthetic consumption experiences is precisely to serve as the primary way to renew the individual's store of non-instrumental content potentially mobilizable in interaction. This content is the main input that goes into status-segregated conversational rituals \citep[125]{collins75}, justifying their existence through seemingly trivial forms of talk: as Kadushin \citeyearpar[797]{kadushin66} noted in his classic study on the subject, urban social circles are produced and reproduced through shared \emph{cultural tastes} and practices. The common knowledge gained through culture talk serves as a cognitive and emotional basis of solidarity within the micro-ritual generated in the interaction, as persons become entranced by ``the joy of sharing names'' \citep{dimaggio87}. This may lead, under typical circumstances, to the confirmation and reification of group attitudes \citep[461]{fine77}. 
